\documentclass[10pt]{article}

\usepackage[margin=24mm]{geometry}
\usepackage{amsmath,amssymb}
\usepackage{booktabs}
\usepackage{xcolor}
\usepackage{hyperref}
\usepackage{fancyhdr}

\definecolor{linkblue}{HTML}{1F5A7A}
\definecolor{softblue}{HTML}{EAF3F7}
\hypersetup{
  colorlinks=true,
  linkcolor=linkblue,
  urlcolor=linkblue,
  citecolor=linkblue,
  pdftitle={Why the Brownian Rate Estimate Is an Average Outer Product},
  pdfauthor={Embryogenesis Pareto analysis project}
}

\setlength{\parindent}{0pt}
\setlength{\parskip}{5pt}
\renewcommand{\arraystretch}{1.15}
\pagestyle{fancy}
\fancyhf{}
\lhead{Brownian covariance MLE}
\rhead{Equation 2 derivation}
\cfoot{\thepage}
\setlength{\headheight}{14pt}

\newcommand{\Normal}{\mathcal{N}}
\newcommand{\Var}{\operatorname{Var}}
\newcommand{\tr}{\operatorname{tr}}

\title{\textbf{Why the Brownian Rate Estimate Is an\\Average Outer Product}\\[2mm]
\large An intuitive and mathematical derivation of Equation 2 in the companion methods note}
\author{Embryogenesis Pareto analysis project}
\date{1 September 2026}

\begin{document}
\maketitle

\begin{abstract}
Equation 2 of the companion parametric-bootstrap methods note states that, for
a Brownian process observed on every parent--child edge, the branch-standardized
increments are independent multivariate-normal observations.  Their
maximum-likelihood covariance is their average outer product.  This note
derives that result first in one dimension, then geometrically and algebraically
in multiple dimensions, and finally reconnects it to branch durations in the
full \textit{C. elegans} lineage analysis.
\end{abstract}

\section{The result in plain language}

For edge $e$, let
\[
  \boldsymbol\Delta_e
  = \mathbf{x}_{child(e)}-\mathbf{x}_{parent(e)}
\]
be its observed 23-dimensional state change, and let $t_e>0$ be the duration of
the edge.  Brownian motion assumes
\[
  \boldsymbol\Delta_e
  \sim \Normal\!\left(\mathbf 0,t_e\boldsymbol\Sigma\right).
\]
Longer edges have proportionally more variance.  Dividing by the square root
of time removes that expected scaling:
\[
  \mathbf z_e=\frac{\boldsymbol\Delta_e}{\sqrt{t_e}}
  \sim\Normal(\mathbf0,\boldsymbol\Sigma).
\]

The 1,000 vectors $\mathbf z_e$ are therefore 1,000 observations of the same
zero-centered Gaussian cloud.  Maximum likelihood chooses the covariance
matrix whose cloud has the same width, orientation, and elongation as the
observed scatter.  Algebraically, that covariance is
\setcounter{equation}{1}
\begin{equation}
  \widehat{\boldsymbol\Sigma}
  =\frac{1}{E}\sum_{e=1}^{E}\mathbf z_e\mathbf z_e^\top
  =\frac{1}{E}\sum_{e=1}^{E}
    \frac{\boldsymbol\Delta_e\boldsymbol\Delta_e^\top}{t_e}.
  \label{eq:target}
\end{equation}

The rest of the note explains why.

\section{Begin with one dimension}

Suppose an edge has only one standardized feature:
\[
  z_e\sim\Normal(0,\sigma^2).
\]
The density of one observation is
\[
  p(z_e\mid\sigma^2)
  =\frac{1}{\sqrt{2\pi\sigma^2}}
   \exp\!\left(-\frac{z_e^2}{2\sigma^2}\right).
\]

There are two competing penalties:

\begin{itemize}
  \item If $\sigma^2$ is too small, observed changes far from zero become very
        unlikely because of the exponential term.
  \item If $\sigma^2$ is too large, probability is spread over an unnecessarily
        wide interval because of the $1/\sqrt{\sigma^2}$ term.
\end{itemize}

The maximum-likelihood variance balances those penalties.

For $E$ independent edges, the likelihood is the product of their densities.
Taking logarithms turns the product into a sum:
\begin{align*}
  \ell(\sigma^2)
  &=\log L(\sigma^2)\\
  &=-\frac{E}{2}\log(2\pi)
    -\frac{E}{2}\log(\sigma^2)
    -\frac{1}{2\sigma^2}\sum_{e=1}^{E}z_e^2.
\end{align*}
The first term is constant with respect to $\sigma^2$.  Differentiating the
other terms gives
\[
  \frac{d\ell}{d\sigma^2}
  =-\frac{E}{2\sigma^2}
   +\frac{1}{2(\sigma^2)^2}\sum_e z_e^2.
\]
At the maximum this derivative is zero, so
\[
 -E\sigma^2+\sum_e z_e^2=0,
 \qquad\text{and hence}\qquad
 \boxed{\widehat{\sigma}^2=\frac{1}{E}\sum_e z_e^2.}
\]

Thus the best-fitting zero-centered Gaussian variance is the average squared
distance from zero.  Small observations produce a narrow fitted Gaussian;
large observations produce a broad one.

\section{The multivariate version: outer products}

For $p$ features, one standardized edge change is a column vector
\[
  \mathbf z_e=(z_{e1},\ldots,z_{ep})^\top.
\]
Instead of a single variance, the Gaussian requires a covariance matrix
\[
  \boldsymbol\Sigma=
  \begin{pmatrix}
    \Var(z_1) & \operatorname{Cov}(z_1,z_2) & \cdots\\
    \operatorname{Cov}(z_2,z_1) & \Var(z_2) & \cdots\\
    \vdots&\vdots&\ddots
  \end{pmatrix}.
\]

The multivariate analogue of squaring $z_e$ is taking its \emph{outer
product}.  In two dimensions,
\[
  \mathbf z_e=\begin{pmatrix}a\\b\end{pmatrix}
  \quad\Longrightarrow\quad
  \mathbf z_e\mathbf z_e^\top
  =\begin{pmatrix}a^2&ab\\ab&b^2\end{pmatrix}.
\]

Each element has a direct interpretation:

\begin{itemize}
  \item $a^2$ records how much feature 1 varied on this edge;
  \item $b^2$ records how much feature 2 varied;
  \item $ab$ records whether they changed together.  It is positive when they
        have the same sign and negative when they have opposite signs.
\end{itemize}

Averaging these matrices over edges therefore averages the observed variances
and covariances simultaneously:
\[
  \widehat{\boldsymbol\Sigma}
  =\frac{1}{E}\sum_e\mathbf z_e\mathbf z_e^\top.
\]

Geometrically, a multivariate Gaussian forms an ellipse or ellipsoid.  The
diagonal entries determine its spread along the feature axes; off-diagonal
entries rotate it toward directions in which features tend to change together.
The average outer product is precisely the matrix describing the empirical
ellipsoid of the observations around the known mean zero.

\subsection{A concrete two-dimensional example}

Consider three standardized edge changes:
\[
  \mathbf z_1=\begin{pmatrix}1\\2\end{pmatrix},\qquad
  \mathbf z_2=\begin{pmatrix}-1\\0\end{pmatrix},\qquad
  \mathbf z_3=\begin{pmatrix}0\\-2\end{pmatrix}.
\]
Their outer products are
\[
\begin{pmatrix}1&2\\2&4\end{pmatrix},\qquad
\begin{pmatrix}1&0\\0&0\end{pmatrix},\qquad
\begin{pmatrix}0&0\\0&4\end{pmatrix}.
\]
Their average is
\[
  \widehat{\boldsymbol\Sigma}
  =\frac{1}{3}\begin{pmatrix}2&2\\2&8\end{pmatrix}
  =\begin{pmatrix}0.67&0.67\\0.67&2.67\end{pmatrix}.
\]
Feature 2 has much greater variance than feature 1, and the positive
off-diagonal entry indicates a positive association.  The fitted Gaussian is
therefore elongated mainly in feature 2 and tilted toward changes in which the
two features move together.

\section{Why the outer-product average maximizes likelihood}

The multivariate-normal log-likelihood is
\begin{equation}
  \ell(\boldsymbol\Sigma)
  =-\frac{E}{2}\log|\boldsymbol\Sigma|
   -\frac{1}{2}\sum_e
      \mathbf z_e^\top\boldsymbol\Sigma^{-1}\mathbf z_e
   +\text{constant}.
  \label{eq:loglik}
\end{equation}

The first term penalizes an unnecessarily large Gaussian volume because
$|\boldsymbol\Sigma|$ measures the squared volume scale of its covariance
ellipsoid.  The second term penalizes a covariance that is too small or points
in the wrong direction.  Its quadratic form is the squared Mahalanobis
distance of an observation from zero under the proposed covariance.

Define the empirical scatter matrix
\[
  \mathbf S=\frac{1}{E}\sum_e\mathbf z_e\mathbf z_e^\top.
\]
Using
\[
  \mathbf z_e^\top\boldsymbol\Sigma^{-1}\mathbf z_e
  =\tr\!\left(\boldsymbol\Sigma^{-1}
               \mathbf z_e\mathbf z_e^\top\right),
\]
Equation~\eqref{eq:loglik} becomes
\[
  \ell(\boldsymbol\Sigma)
  =-\frac{E}{2}\left[
      \log|\boldsymbol\Sigma|
      +\tr(\boldsymbol\Sigma^{-1}\mathbf S)
    \right]+\text{constant}.
\]

The relevant matrix derivatives are
\[
 \frac{\partial\log|\boldsymbol\Sigma|}{\partial\boldsymbol\Sigma}
 =\boldsymbol\Sigma^{-1},
 \qquad
 \frac{\partial\tr(\boldsymbol\Sigma^{-1}\mathbf S)}
      {\partial\boldsymbol\Sigma}
 =-\boldsymbol\Sigma^{-1}\mathbf S\boldsymbol\Sigma^{-1}.
\]
Setting the derivative of the log-likelihood to zero gives
\[
 -\boldsymbol\Sigma^{-1}
 +\boldsymbol\Sigma^{-1}\mathbf S\boldsymbol\Sigma^{-1}=0.
\]
Multiplying on both sides by $\boldsymbol\Sigma$ yields
\[
  \boxed{\boldsymbol\Sigma=\mathbf S.}
\]
Therefore the Gaussian whose covariance ellipsoid matches the empirical
scatter is not only intuitive: it is exactly the maximum-likelihood solution.

\section{Returning to branch durations}

Because
\[
  \mathbf z_e=\frac{\boldsymbol\Delta_e}{\sqrt{t_e}},
\]
its outer product is
\[
  \mathbf z_e\mathbf z_e^\top
  =\frac{\boldsymbol\Delta_e\boldsymbol\Delta_e^\top}{t_e}.
\]
Substitution into the average scatter matrix recovers Equation~\eqref{eq:target}:
\[
  \boxed{
  \widehat{\boldsymbol\Sigma}
  =\frac{1}{E}\sum_e
    \frac{\boldsymbol\Delta_e\boldsymbol\Delta_e^\top}{t_e}.}
\]

In words: compute each edge's multivariate squared change per unit time,
including how features change together, and average that information over all
edges.

\section{Why divide by $E$ rather than $E-1$?}

The Brownian model fixes the expected increment at zero.  Because this mean is
known rather than estimated from the edge increments, no degree of freedom is
spent estimating it.  Consequently, division by $E$ is both the
maximum-likelihood estimator and an unbiased estimator under the stated model.

If a nonzero drift vector were estimated from the same increments, the
maximum-likelihood covariance would still use an $E$ denominator, but a
degrees-of-freedom-adjusted covariance estimator would be required for
unbiasedness.  This is the same distinction as the familiar $n$ versus $n-1$
denominator in elementary statistics.

\section{Direct phylogenetic connection and assumptions}

Revell and Collar derive the corresponding multivariate Brownian rate-matrix
MLE for correlated tip observations \cite{revellcollar2009}:
\[
  \widehat{\mathbf R}
  =\frac{(\mathbf X-\mathbf1\mathbf a^\top)^\top
          \mathbf C^{-1}
         (\mathbf X-\mathbf1\mathbf a^\top)}{n},
\]
where $\mathbf C$ is the covariance implied by shared ancestry.  In our
fully-observed-edge setting, branch standardization transforms the data into
independent increments, so the covariance structure becomes the identity and
their expression reduces to $\mathbf Z^\top\mathbf Z/E$.  Felsenstein's
independent-contrast construction provides the foundational connection between
Brownian motion, branch-length standardization, and independent Gaussian
changes \cite{felsenstein1985}.  Freckleton gives a clear likelihood treatment
of Brownian rate estimation and the ML versus degrees-of-freedom denominators
\cite{freckleton2012}.

Equation~\eqref{eq:target} relies on the following assumptions:

\begin{itemize}
  \item the expected increment is zero unless drift is modeled explicitly;
  \item Brownian increments are independent across edges conditional on their
        observed parent states;
  \item increment covariance is proportional to branch duration;
  \item one common $\boldsymbol\Sigma$ applies to all scored edges; and
  \item node states are observed without substantial measurement error.
\end{itemize}

The last condition matters.  If the observed state is
$\mathbf y_v=\mathbf x_v+\boldsymbol\eta_v$, adjacent observed increments share
the measurement error of their common node.  They are then correlated even if
the latent Brownian increments are independent.  A publication analysis should
therefore either justify negligible measurement error or extend the model to
include an observation-error covariance.

\begin{thebibliography}{9}

\bibitem{felsenstein1985}
Felsenstein, J. (1985).
Phylogenies and the comparative method.
\textit{The American Naturalist}, 125(1), 1--15.
\href{https://doi.org/10.1086/284325}{doi:10.1086/284325}.

\bibitem{revellcollar2009}
Revell, L. J., and Collar, D. C. (2009).
Phylogenetic analysis of the evolutionary correlation using likelihood.
\textit{Evolution}, 63(4), 1090--1100.
\href{https://doi.org/10.1111/j.1558-5646.2009.00616.x}%
{doi:10.1111/j.1558-5646.2009.00616.x}.

\bibitem{freckleton2012}
Freckleton, R. P. (2012).
Fast likelihood calculations for comparative analyses.
\textit{Methods in Ecology and Evolution}, 3(5), 940--947.
\href{https://doi.org/10.1111/j.2041-210X.2012.00220.x}%
{doi:10.1111/j.2041-210X.2012.00220.x}.

\end{thebibliography}

\end{document}
